\documentclass[aspectratio=169]{beamer}

\usetheme{Madrid}
\usecolortheme{default}

\usepackage{booktabs}
\usepackage{hyperref}
\usepackage{xcolor}
\usepackage{array}
\usepackage{graphicx}

% Global bullet spacing
\newlength{\myitemsep}
\setlength{\myitemsep}{9pt}
\let\olditemize\itemize
\renewcommand{\itemize}{\olditemize\setlength{\itemsep}{\myitemsep}}
\let\oldenumerate\enumerate
\renewcommand{\enumerate}{\oldenumerate\setlength{\itemsep}{\myitemsep}}

% ============================================================
%  HOUSE STYLE (agreed with Bruno, 2 Sept 2026)
%    - No em-dashes. Use colons.
%    - No boldface anywhere.
%    - No loose "punchline" sentence at the foot of a slide.
%    - American spelling.
%    - Citations live on the References slide at the end.
% ============================================================

\title{The State of AI}
\subtitle{D-Lab AI Pulse, Fall 2026: Session 1}
\author[D-Lab]{Materials: \url{https://github.com/dlab-berkeley/AI_Pulse}\\Workshop suggestions: \url{https://forms.gle/x8KfejgCViE7t6iQA}}
\date{1 September 2026}

\begin{document}

\begin{frame}
    \titlepage
\end{frame}

\begin{frame}{Today's Roadmap}
    \begin{itemize}\setlength{\itemsep}{7pt}
        \item AI Pulse is D-Lab's bi-weekly series on AI tools for academia, no experience required!
        \item Part 1, where things stand, what changed, what exists now, and what it costs
        \item Part 2, what agents can do, and are these tools actually working?
        \item Part 3, how your role changes, what to hand over, and what that leaves you doing
    \end{itemize}
\end{frame}

\begin{frame}{This Semester}
    Eight sessions, 1 September to 8 December. Every other Tuesday.

    \vspace{1em}
    \begin{itemize}
        \item Workshops alternate: ``wide-view'' sessions, and hands-on sessions with demos
        \item Office hours on the weeks in between, same day and time, on Zoom
        \item Office hours are for your actual problem: getting a tool working, or working out
              which one fits what you are trying to do
        \item Coming up: The Major AI Ecosystems (15 Sept) and AI in the News (29 Sept)
        \item Materials: \url{https://github.com/dlab-berkeley/AI_Pulse} $\cdot$
              D-Lab: \url{https://dlab.berkeley.edu}
    \end{itemize}
\end{frame}

% ============================================================
%  PART 1
% ============================================================

\begin{frame}[plain]
    \vfill
    \begin{center}
        {\Huge Part 1}\\[0.8em]
        {\Large Where Things Stand}
    \end{center}
    \vfill
\end{frame}

% Gowers quote verified: Timothy Gowers on the Erdos unit distance proof,
% "I would have recommended acceptance without any hesitation."
% Unit distance conjecture posed 1946, disproved 20 May 2026 by an internal OpenAI
% general-purpose reasoning model. Constructive proof, externally verified.
\begin{frame}{How We Got Here}
    \begin{itemize}\setlength{\itemsep}{7pt}
        \item 2017: \textit{Attention Is All You Need}, the architecture behind all of it
        \item 2020: AlphaFold 2 solves protein folding, earning its authors the 2024
              Nobel Prize in Chemistry
        \item 2022: ChatGPT opens to the public, five years after the paper
        \item 2025: Claude Code. AI goes from chats to agents
        \item May 2026: the first AI-produced proof of a major open problem. An 80-year-old
              Erd\H{o}s conjecture, reasoned through rather than searched for
    \end{itemize}

    \vspace{0.9em}
    Timothy Gowers, Fields Medallist: \textit{``if a human had written the paper and
    submitted it to the Annals of Mathematics \ldots I would have recommended acceptance
    without any hesitation.''}
\end{frame}

\begin{frame}{From Chat to Agents}
    \begin{itemize}
        \item 2023, chats
        \begin{itemize}
            \item A powered-up search engine. You asked; you copied the answer out by hand
        \end{itemize}
        \item 2024, files
        \begin{itemize}
            \item It could read your documents, and hand one back finished
        \end{itemize}
        \item 2025, tasks
        \begin{itemize}
            \item You describe the job, walk away, come back to the result
        \end{itemize}
        \item 2026, agents
        \begin{itemize}
            \item A whole team. Your job becomes that of a manager
        \end{itemize}
    \end{itemize}

    \vspace{1.2em}
Each step changed what you can ask for, not just how much you can hand over.
\end{frame}

% Boris Cherny X thread 2 Jan 2026; Sequoia AI Ascent 2026; Lenny's Podcast.
% 80% figure: Anthropic, code merged to production May 2026 (VentureBeat, 4 Jun 2026).
% 90% and 100% also circulate depending on what is counted.
\begin{frame}{Fifteen Claudes at Once}
    Boris Cherny, who built Claude Code, on a normal working day:

    \vspace{0.8em}
    \begin{itemize}
        \item Fifteen running at once: five in terminals, five to ten in the browser,
              one on his phone
        \item Each on its own task, in its own copy of the project, so nothing collides
        \item They test their own work before handing it back; he reviews and decides what
              is kept
        \item Ten to thirty finished pieces of work merged per day
    \end{itemize}

    \vspace{1.2em}
    At Anthropic, 80\% of the code merged into production in May 2026 was written by Claude.
\end{frame}

% Gemini: Google, 11 Aug 2026, 1B monthly active users.
% Gallup Q1 2026: 50% of US employees used AI at work, up from 46%; 28% weekly.
% Pangram, 18 Nov 2025: 21% of ICLR 2026's 70,000 reviews fully AI-generated.
% HEPI/Kortext 2026: 95% of 1,054 UK undergraduates; 94% on assessed work.
% PwC Global AI Jobs Barometer 2026: AI-specialist postings +68.9% vs +8.6% overall.
\begin{frame}{Meanwhile, Everybody Else}
    \begin{itemize}
        \item One billion people used Google's Gemini app last month, the fastest growth of
              any product Google has shipped. ChatGPT is in the same range
        \item Half of US employees used AI at work this year, up from 46\% a year ago.
              More than a quarter use it every week
        \item At ICLR 2026, the largest conference in machine learning, 21\% of the 70,000
              peer reviews were fully AI-written, and over half showed some AI
        \item 95\% of UK undergraduates use AI, 94\% of them on assessed work
        \item Job ads asking for AI skills grew 69\% in a year, while hiring overall grew 9\%
    \end{itemize}

    \vspace{1em}
    Is it a way of standing out, or is it necessary not to lag behind?
\end{frame}

% Photograph: 1960s drafting hall. Rights holder unidentified; gitignored, not republished.
\begin{frame}{How Is Work Changing?}
    \begin{columns}[c]
    \begin{column}{0.44\textwidth}
        \begin{center}
        \IfFileExists{drafting_room.jpg}{%
            \includegraphics[width=\linewidth]{drafting_room.jpg}%
        }{%
            \small\itshape [Photograph omitted from the public repository:\\
            a 1960s drafting hall, rights holder unidentified.]%
        }
        \end{center}
    \end{column}
    \begin{column}{0.52\textwidth}
        \begin{itemize}\setlength{\itemsep}{10pt}
            \item Did engineers disappear when CAD was released?
            \item Could engineers keep their jobs without mastering the new tech?
            \item How long did it take for CAD to become necessary, and how long will it
                  take for AI?
        \end{itemize}
    \end{column}
    \end{columns}
\end{frame}

% Fastest-rotting slide in the deck. Re-check every product name before each delivery.
% Devin is current: Cognition raised $1B at a $26B valuation, June 2026.
\begin{frame}{What Is Actually Out There}
    \begin{itemize}\setlength{\itemsep}{4pt}
        \item General assistants: ChatGPT, Claude, Gemini, Grok
        \vspace{2pt}\begin{itemize}\item They all do the same things now. Choose on fit, not on capability\end{itemize}
        \item Coding platforms: Cursor, Claude Code, Devin, Copilot
        \vspace{2pt}\begin{itemize}\item Built for code, but increasingly just run tasks. A platform, not a chatbot\end{itemize}
        \item Research tools: NotebookLM, Elicit, Consensus
        \vspace{2pt}\begin{itemize}\item Finding, reading and synthesizing literature, with the sources attached\end{itemize}
        \item Generation: Midjourney, Veo, ElevenLabs, Lovable
        \vspace{2pt}\begin{itemize}\item Images, video, voice, websites. One job each, done well\end{itemize}
        \item Productivity: Granola, Otter, Notion
        \vspace{2pt}\begin{itemize}\item Ready-made. They take the grunt work without you designing anything\end{itemize}
    \end{itemize}
\end{frame}

% Berkeley: Gemini, NotebookLM, Zoom AI Companion covered by campus agreements.
% GEMINI CLI: Google's docs say the 1,000 req/day free tier needs a PERSONAL Google account
% and that Workspace/EDU accounts need a paid Code Assist licence. Bruno used it on his
% Berkeley account until ~May 2026. Hence the asterisk.
% Plan contents: OpenAI per chatgpt.com/pricing; Anthropic per claude.com/pricing.
\begin{frame}{What It Costs}
    \begin{itemize}\setlength{\itemsep}{6pt}
        \item Free with your Berkeley account: Gemini, NotebookLM and Gemini CLI$^{*}$
        \begin{itemize}
            \item Under the campus agreement your data does not train the model, and it
                  reaches your Drive, Docs, Gmail and Calendar
        \end{itemize}
        \item OpenAI: Free $\cdot$ Go \$8 $\cdot$ Plus \$20 $\cdot$ Pro \$100 or \$200 a month
        \begin{itemize}
            \item Free is chat; Go adds messages, uploads and voice; Plus adds deep research
                  and projects; Pro adds the reasoning models and research previews
        \end{itemize}
        \item Anthropic: Free $\cdot$ Pro \$20 $\cdot$ Max \$100 or \$200 a month
        \begin{itemize}
            \item Free is chat; Pro adds Claude Code, Cowork and the Microsoft 365
                  integration; Max adds usage limits and priority access
        \end{itemize}
    \end{itemize}

    \vspace{0.8em}
    \begin{center}
    Next session: OpenAI vs Anthropic in depth, with demos.
    \end{center}

    \vspace{0.4em}
    {\small\itshape $^{*}$ Evolving situation.}
\end{frame}

% Prices verified 2 Sept 2026. Anthropic from platform.claude.com/docs; OpenAI from
% CloudZero's table, updated 20 Aug 2026, after the 30 July cut (Terra -20%, Luna -80%).
% Output costs ~5x input in every row, both vendors.
% Table is a deliberate exception to the no-boxes rule: six prices need columns.
\begin{frame}{What a Token Costs}
    \begin{itemize}\setlength{\itemsep}{5pt}
        \item A token is the fundamental unit of information for AI: roughly three quarters
              of a word
        \item You pay for what you send in, and again for what comes back
        \item Output costs roughly five times input, in every model, from both vendors
    \end{itemize}

    \vspace{0.6em}
    \begin{center}
    \small
    \begin{tabular}{lrr@{\hskip 2.2em}lrr}
        \toprule
        \multicolumn{3}{c}{Anthropic} & \multicolumn{3}{c}{OpenAI} \\
        \cmidrule(r){1-3}\cmidrule(l){4-6}
        Model & In & Out & Model & In & Out \\
        \midrule
        Fable 5     & \$10 & \$50 & GPT-5.6 Sol   & \$5    & \$30   \\
        Opus 5      & \$5  & \$25 & GPT-5.6 Terra & \$2    & \$12   \\
        Sonnet 5    & \$2  & \$10 & GPT-5.6 Luna  & \$0.20 & \$1.20 \\
        Haiku 4.5   & \$1  & \$5  &               &        &        \\
        \bottomrule
    \end{tabular}
    \\[0.3em]
    {\scriptsize US dollars per million tokens. September 2026.}
    \end{center}
\end{frame}

% LMArena, September 2026: blind pairwise voting. Top five within 20 Elo points;
% #1 to #10 is 32 points. #1 is Claude Opus 4.6 (Fast) at 1500, AHEAD of the newer
% Fable 5 at 1494. Gemini 3.5 Flash, a cheap fast model, is #4 at 1480.
% Fable 5 input is 50x Luna's ($10 vs $0.20).
\begin{frame}{Older Models Are Good Enough}
    \begin{itemize}\setlength{\itemsep}{7pt}
        \item On LMArena, where people vote blind on which of two answers they prefer, the
              top five models sit within 10 points of each other
        \item Second place goes to Claude Opus 4.6, a model two generations old, two points
              behind the leader and at half the price
        \item Gemini 3.7 Flash sits ninth, 17 points off the top, at a thirteenth of the
              leader's input price
        \item What decides is interface, price, privacy rules, and what your institution
              already pays for
        \item We have reached a stage in which the best option is not the best model
    \end{itemize}
\end{frame}

% ============================================================
%  PART 2
% ============================================================

\begin{frame}[plain]
    \vfill
    \begin{center}
        {\Huge Part 2}\\[0.8em]
        {\Large What Agents Do, and Whether It Works}
    \end{center}
    \vfill
\end{frame}

\begin{frame}{What Makes It an Agent}
    \begin{itemize}\setlength{\itemsep}{7pt}
        \item A chat produces text. You are still the one who acts on it
        \item An agent can open, read, write, run and check things itself
        \item The model underneath is the same. What is new is access and permission,
              not intelligence
        \item It started with code because code checks itself: it runs or it does not,
              and the tests pass or they fail
        \item So the model can try, fail and fix before a human sees anything, and keep
              going until the task is done or it needs you
    \end{itemize}
\end{frame}

\begin{frame}{What an Agent Can Reach}
    \begin{itemize}\setlength{\itemsep}{6pt}
        \item Your files: read them, edit them, write new ones
        \item Your calendar, mail and notes, through MCP, the common standard that lets tools
              connect to each other
        \item Code: write it, run it, read the error, fix it, run it again
        \item The internet: search it, open pages, and use what it finds
        \item Your team's tools: GitHub, Slack, Jira and the like, through that same standard
        \item Time and place: start a task from your phone, or schedule it to run every day
    \end{itemize}
\end{frame}

% Gateways: OpenRouter, LiteLLM. One interface, many models, no rewriting to switch.
\begin{frame}{Beyond the Main AI Apps}
    \begin{itemize}\setlength{\itemsep}{6pt}
        \item Context is the constraint for repeated, independent tasks: an agent carries one
              conversation, and a thousand items need a thousand clean starts
        \item So for the same operation across thousands of items, you call the model directly
              through its API instead
        \item A gateway such as OpenRouter or LiteLLM sends that call to any model, and lets
              you switch models without rewriting anything
        \item Four thousand PDFs, one extraction schema, one spreadsheet of results
        \item You can also adapt an open model such as Llama to one narrow task: sensitive
              data stays on your machine, cost per document falls, and the model can be
              specialized
    \end{itemize}
\end{frame}

% Elicit: Cambridge "Research Synthesis Methods" feasibility study on its data extraction.
% SciSpace (formerly Typeset): 280M+ papers indexed.
% AI Co-Scientist: Google DeepMind, Nature paper May 2026, wet-lab validated, not clinical.
% MHS: Anthropic research preview, announced Thursday 27 Aug 2026. CMU integrated in 8 hours;
% UW connected six instruments in under a week. Selected partners only.
\begin{frame}{Built for Science}
    \begin{itemize}\setlength{\itemsep}{7pt}
        \item Elicit: works across the literature, searching, screening and extracting at
              systematic-review scale
        \item SciSpace: sits inside one paper with you. Chat with the PDF, get a figure or a
              derivation explained, pull the numbers out
        \item AI Co-Scientist: proposes hypotheses and designs the experiments to test them.
              Published in \textit{Nature} in May, validated at the bench
    \end{itemize}

    \vspace{0.9em}
    \begin{center}
    Two sessions on this later in the semester: AI in Science, Workflows, and AI in Science,
    Tools.
    \end{center}

    \vspace{0.7em}
    {\small \textbf{Developing:} on 27 August Anthropic released the Model Hardware
    Standard, which is MCP for physical instruments. Robotic arms, lab equipment,
    manufacturing machines.}
\end{frame}

% Kimi K3: Moonshot AI, released 16 July 2026, weights 27 July. 2.8T parameters, the largest
% open-weight model released to date. #1 on Frontend Code Arena (1,679); 3rd on GDPval-AA v2.
% Open weights on permissive licences: DeepSeek, Qwen, GLM, MiniMax remain dramatically
% cheaper. Airbnb was questioned by Congress in April 2026 over using Qwen; a self-hosted
% open-weight model sends nothing to its author.
\begin{frame}{Nobody Stays in Front}
    \begin{itemize}\setlength{\itemsep}{7pt}
        \item The lead changes every few months. Whoever tops the table today has usually
              been passed by the next session of this workshop
        \item In July, Moonshot AI released Kimi K3: 2.8 trillion parameters, the largest
              openly released model to date. It placed first on one major coding leaderboard
              and third on another
        \item Chinese labs are now at the frontier, and the open-weight ones,
              DeepSeek, Qwen and GLM among them, remain far cheaper than the American models
        \item Open weights matter more than headline price: a model you run yourself costs
              nothing per use, and sends nothing anywhere
        \item So weigh the cost of switching against the gain from always being on the best
    \end{itemize}
\end{frame}


\begin{frame}{Changing the Question}
    \begin{itemize}\setlength{\itemsep}{8pt}
        \item Everything so far has been about what these tools can do, and most of it comes
              from the people who make them
        \item Very little of it comes from a study designed to find out whether the tools
              actually help
        \item So the question changes: when somebody measures carefully, what do they find?
        \item Three studies: a state labor department, a hospital, and a peer review process
    \end{itemize}
\end{frame}

% Stanford RegLab with the Colorado Dept. of Labor and the US DOL, published January 2026.
% Fine-tuned LLaMA-3 in a secure government sandbox. RCT: 8 adjudicators randomized at week
% level, 788 drafts, blind QA against two control arms.
% Time saved: 4 seconds, p = 0.8. Quality vs adjudicator alone: no difference, p = 0.37.
\begin{frame}{What Happens When Somebody Measures}
    \begin{itemize}\setlength{\itemsep}{6pt}
        \item When you file for unemployment insurance, an adjudicator reviews your claim and
              writes back asking for whatever is missing. It is slow, and it is repetitive
        \item Colorado's labor department and Stanford fine-tuned a model on the department's
              own records to draft those follow-up questions
        \item They ran a randomized trial: 788 drafts, eight adjudicators, blind quality
              review against two control groups
        \item Time saved per case: four seconds. Statistically indistinguishable from zero
        \item Quality, against the adjudicator working alone: no significant difference.
              They published the result and did not deploy the system
    \end{itemize}
\end{frame}

% NOHARM: arXiv 2512.01241, Goh et al., 46 co-authors, v4 13 Jul 2026. 1,100 consultation
% cases, 10 specialties, 20 LLMs + 4 clinical AI tools, plus an RCT of 101 physicians.
% Severe-harm potential in up to 24.6% of cases; >80% of severe errors were OMISSIONS.
% UK Government "Consult": reviewers used "Other" 793 times against the model's 47, a 17x gap.
\begin{frame}{The Failure You Cannot See}
    \begin{itemize}\setlength{\itemsep}{7pt}
        \item A model does not know what it does not know
        \item In a 1,100-case medical safety study, over 80\% of severe errors were things
              the model left out, not things it invented
        \item A UK government review of public consultation responses found the same shape:
              human reviewers used the ``Other'' category seventeen times more often than the
              model did
        \item It does not say that it is unsure. It files the case under the nearest heading
              it already has, and nothing in the output shows that anything is missing
    \end{itemize}
\end{frame}

% Harvard/Stanford, Science: 76 ER cases, o1-preview only, diagnostic accuracy.
% NOHARM: 1,100 consultations, 10 specialties, 24 systems, safety.
\begin{frame}{Accuracy Is Not Safety}
    \begin{itemize}\setlength{\itemsep}{7pt}
        \item Two studies of AI in clinical medicine, months apart, reached opposite
              headlines. Both are correct
        \item Harvard and Stanford: 76 emergency cases, one model. It matched or exceeded
              expert physicians on diagnosis
        \item NOHARM: 1,100 consultations, 24 systems. Potential for severe harm in up to a
              quarter of cases
        \item One measured whether it gets the answer right. The other measured what happens
              when it does not
        \item The average and the tail are different questions, and only one of them made the
              headlines
    \end{itemize}
\end{frame}

% ============================================================
%  PART 3
% ============================================================

\begin{frame}[plain]
    \vfill
    \begin{center}
        {\Huge Part 3}\\[0.8em]
        {\Large How Your Role Changes}
    \end{center}
    \vfill
\end{frame}

\begin{frame}{Replacing Workers, or Replacing Work?}
    \begin{itemize}\setlength{\itemsep}{8pt}
        \item The question people ask is whether these tools replace people
        \item A more immediate question is: which parts of a job they take over, and which
              parts they leave alone
        \item Every technology that automated a task changed what the job was long before it
              changed how many jobs there were
        \item So the question is not whether your field survives this, but which hours of
              your week are the ones under discussion
    \end{itemize}
\end{frame}

% Kasparov lost the Deep Blue rematch on 11 May 1997, 3.5 to 2.5.
% PAL/CSS Freestyle tournament, 2005. Winners: Steven Cramton and Zackary Stephen ("ZackS").
% Quote from "The Chess Master and the Computer", NYRB, 11 Feb 2010.
\begin{frame}{The Centaur}
    \begin{itemize}\setlength{\itemsep}{6pt}
        \item In May 1997, Deep Blue beat Garry Kasparov. Chess was widely treated as settled
        \item Kasparov's response was to invent a format where humans and machines play
              together. Eight years later a tournament ran it open to anyone, with any machine
        \item It was won by two American amateurs running three ordinary PCs
        \item They were not better at chess. They were better at running the machines:
              knowing which positions to push on, and when to disagree
        \item Kasparov: \textit{``weak human + machine + better process was superior to a
              strong computer alone''}, and also to a strong human with a worse process
    \end{itemize}
\end{frame}

% VERIFIED 2 Sep 2026 by reading erdosproblems.com and Tao's wiki directly in the browser.
% Site front page: 1,220 problems, 576 (47%) solved.
% Pre-AI rate: Quanta, 111 problems moved open->solved between 2024 and Aug 2025
%   = 20 months = ~5.5/month.
% D-DAY: Erdos Problem 90 (unit distance), OpenAI internal model, 20 MAY 2026, full solution.
%   Claude Mythos solved it independently on 26 May 2026. Both in wiki section 1(a).
% THE 18: counted by hand from Tao's wiki section 1(a), "AI standalone" (human involvement
% non-significant), taking only entries marked green/full solution and counting DISTINCT
% problems. They are: 38, 90, 125, 205, 457, 619, 694, 741, 960, 987, 990, 1014, 1051,
% 1091, 1141, 1196, 1202, 1217 = 18. ([90] appears twice: OpenAI 20 May and Claude Mythos
% 26 May, solved independently. Counted once.)
% THE 36: same method on section 1(d), "AI collaborating with humans", green only, distinct
% problems. 33 of the 36 fall in 2026; the other three began in Oct-Dec 2025.
% Both counts run to the wiki's cutoff: it stopped being updated on 30 June 2026.
% !! TIMESPAN: the 18 run from 10 Jan to 9 Jun 2026, i.e. FIVE MONTHS, not six weeks.
%    (13 of the 18 do fall between 1 Apr and 20 May, which is about seven weeks, if a
%    tighter window is ever wanted. The slide deliberately makes no time claim.)
% !! 16 OF THE 18 PREDATE D-DAY. Before 20 May: 205, 1051, 457, 125, 1202, 960, 987, 990,
%    1091, 1141, 1196, 741, 1217, 1014, 38, 694. On or after: only 90 (20 May) and 619 (9 Jun).
%    So do NOT say "after the unit distance proof, AI solved 18" -- it is the wrong way round.
%    Problem 90 was the first AI proof of a MAJOR problem; 728 (6 Jan) was the first at all.
% LIAM PRICE: Quanta, "Price studied some math in college but left before finishing" --
%    no degree at all. Problem [205], one of the 18 AI-standalone solutions, was the one
%    he elicited (GPT-5.2 Thinking, formalised by Aristotle) -- source: Kevin Barreto's
%    own account on erdosproblems.com, blog:2, 26 Jan 2026.
% Excluded throughout: partial progress, candidate, unverified and incorrect entries.
% NOT from Quanta. The full Quanta text was read in the browser on 2 Sep 2026 and contains
% no figure of 18: its only tallies are 111 (2024-Aug 2025), DeepMind's 4 solved + 9
% rediscovered (Jan), 9 of 353 (May), Astra's 3 more (1 Aug), and 565 solved / 652 open.
% !! CAVEAT to say out loud: the site's own note under OPEN->SOLVED says it is NOT a record
%    of when solutions actually happened or how quickly problems are being solved.
% Quanta, 3 Aug 2026.
% 2026: DeepMind resolved 9 of 353 open problems autonomously (May); a DeepMind team solved
% 4 and found prior solutions to 9 more (January, using Gemini); OpenAI's model settled the
% unit distance conjecture (20 May); OpenAI's Astra made 10 advances including 3 more Erdos
% problems (1 Aug). Barreto (Cambridge undergraduate) and Price (did not finish his degree),
% both early 20s, co-authored a paper with Terence Tao and Jared Duker Lichtman resolving
% Problem 1196. Their earlier Problem 333 turned out to have been solved before.
\begin{frame}{The Same Pattern, in Mathematics}
    \begin{itemize}\setlength{\itemsep}{6pt}
        \item The Erd\H{o}s problems are a public list of open problems in mathematics,
              kept on a website anyone can post to
        \item Professors and hobbyists comment side by side. Terence Tao was among the
              most active
        \item Before AI, the problems closed at about five a month
        \item 18 have now been fully solved by AI, and 36 more with AI assistance
        \item That includes one by Liam Price, who does not have a degree
        \item Knowing how to guide AI was enough
    \end{itemize}
\end{frame}

\begin{frame}{What That Makes Your Job}
    \begin{itemize}\setlength{\itemsep}{7pt}
        \item You are not being asked to write less. You are being asked to specify, review
              and decide more
        \item Describing a task precisely is now a core skill. ``Analyze this'' is not an
              instruction
        \item The tools that coordinate work matter more, not less: version control,
              GitHub, Slack, Monday
        \item You no longer have to memorize the syntax. You do have to be able to tell when
              the answer is wrong
        \item Hand over the work you are time-constrained on, not the work you are
              cognitively constrained on
        \vspace{2pt}
        \begin{itemize}
            \item Ask it to check the grammar across a long document. Do not ask it to write
                  the page for you
        \end{itemize}
    \end{itemize}
\end{frame}

\begin{frame}{Food for Thought}
    \begin{itemize}\setlength{\itemsep}{9pt}
        \item How do I break my work into smaller steps, and what does ``correct'' mean for
              each step?
        \item Centaur chess is over: the engines got strong enough that the human stopped
              adding anything. Which parts of your work are like chess?
        \item What in your work is the thinking, and what is the drafting?
        \item Is using this a way of standing out, or is it keeping up?
        \item If a model wrote the first draft, whose work is it?
    \end{itemize}
\end{frame}

\begin{frame}{Where to Go From Here}
    \begin{itemize}\setlength{\itemsep}{8pt}
        \item Slides and materials: \url{https://github.com/dlab-berkeley/AI_Pulse}
        \item Workshop suggestions: \url{https://forms.gle/x8KfejgCViE7t6iQA}
        \item What Berkeley already gives you:
              \url{https://life.berkeley.edu/ai-tools-for-berkeley-students}
        \item D-Lab: \url{https://dlab.berkeley.edu}
    \end{itemize}

    \vspace{1em}
    \begin{center}
    Next: The Major AI Ecosystems, 15 September.
    Office hours next week, same day and time, on Zoom.
    \end{center}
\end{frame}

\begin{frame}{References}
    \scriptsize
    \begin{itemize}\setlength{\itemsep}{2.5pt}
        \item Gowers quote and the Erd\H{o}s unit distance proof: \textit{Quanta Magazine},
              ``Why the Legendary Erd\H{o}s Problems Are Falling to AI'', 3 August 2026;
              \textit{Nature}, 2026. Problem counts from Bloom's database, \url{erdosproblems.com}
        \item Boris Cherny's workflow: his own thread, 2 January 2026, and Sequoia AI Ascent 2026.
              The 80\% figure: Anthropic, code merged to production May 2026, via VentureBeat,
              4 June 2026
        \item Gemini's billion users: Google, 11 August 2026. Workplace use: Gallup, Q1 2026.
              Student use: HEPI and Kortext, \textit{Student Generative AI Survey 2026},
              1,054 UK undergraduates. Job ads: PwC \textit{Global AI Jobs Barometer 2026}
        \item AI-written peer reviews at ICLR 2026: Pangram, 18 November 2025, over 70,000
              reviews. Note that Pangram sells AI detection
        \item Berkeley's tools and data protections: Berkeley IT and the EVCP office, 2025;
              \url{life.berkeley.edu/ai-tools-for-berkeley-students}
        \item Prices: \url{platform.claude.com/docs} and OpenAI's pricing page, both checked
              September 2026. Leaderboard standings: LMArena, September 2026
        \item Kimi K3: Moonshot AI, 16 July 2026. AI Co-Scientist: Google DeepMind,
              \textit{Nature}, May 2026. Model Hardware Standard: Anthropic, 27 August 2026
        \item The Colorado trial: Stanford RegLab with the Colorado Department of Labor and
              Employment, January 2026. Medical safety: Goh et al., ``First, do NOHARM'',
              arXiv:2512.01241. Diagnosis: Harvard and Stanford, \textit{Science}
        \item Kasparov: ``The Chess Master and the Computer'', \textit{The New York Review of
              Books}, 11 February 2010
    \end{itemize}
\end{frame}

\end{document}
